%% kapa: exam=2 round=30 year=2019 period=1 subject=appraisal_practice subject_ko=감정평가실무 seq=2 number=2 points=30
%% source: Q-Net 감정평가사 기출문제 (cst00309) artlSeq=5205595 file=raw/1008/5205595/감정평가실무.pdf
%% conversion: kordoc 4.24.51 -> tools/gen2.py (auto); review=proofread-sonnet; flags=none
\begin{longquestion}{【문제 2】}
A감정평가법인에 근무 중인 감정평가사 甲은 B도 C시장으로부터 보상목적의 감정평가를 의뢰받아 사전조사 및 현장조사를 마쳤는바, 관련법규 및 이론을 참작하고 제시된 자료를 활용하여 감정평가액을 산출 및 결정하시오. (30점) \par
< 자료 1 > 감정평가 의뢰 내역(요약) \par
1. 의뢰인: B도 C시장 \par
2. 의뢰일자: 2019. 05. 10 \par
3. 가격시점: 2019. 06. 29 \par
4. 공익사업의 명칭: OOO공원 조성사업 \par
5. 의뢰목록(일부 발췌) \par
{\small\centering
\begin{tabular}{|c|c|c|c|c|c|c|}
\hline
\makecell{일련\newline{}번호} & \makecell{소재지} & \makecell{지번} & \makecell{지목(실제)} & \makecell{면적(m\textsuperscript{2})} & \makecell{용도지역} & \makecell{비고} \\
\hline
\makecell{1} & \makecell{B도 C시 D동} & \makecell{148} & \makecell{전(전)} & \makecell{1,235.0} & \makecell{자연녹지} & \makecell{공원 100\%} \\
\hline
\end{tabular}\par}

< 자료 2 > 사업개요 \par
1. 사업계획의 개요 \par
1) 사업명칭: OOO공원 조성사업 \par
2) 사업시행자: B도 C시장 \par
3) 위치: B도 C시 D동 100번지 일원 \par
4) 사업면적: 1,028,520 m\textsuperscript{2}(1단계 462,800 m\textsuperscript{2}, 2단계 565,720 m\textsuperscript{2} 중 1단계 사업) \par
5) 사업기간: 2018. 10. 01∼2020. 12. 31 \par
2. 사업추진 주요 경과 \par
1) 2018. 01. 10: OOO공원 조성계획 결정(변경) 공람공고 \par
2) 2018. 05. 30: OOO공원 조성계획 결정(변경) 및 지형도면 고시 \par
3) 2018. 10. 01: 보상계획 열람 공고 \par
4) 2018. 12. 10: 공익사업 준비를 위한 토지 출입 허가 공고 \par
5) 2019. 05. 10: 감정평가 의뢰 \par
< 자료 3 > 공시지가 표준지, 매매사례 및 평가사례 \par
1. 사업구역 및 인근의 공시지가 표준지 내역 \par
{\small\centering\footnotesize\setlength{\tabcolsep}{2pt}
\begin{tabular}{|c|c|c|c|c|c|c|c|c|c|c|}
\hline
\makecell{기호} & \makecell{소재지} & \makecell{면적\newline{}(m\textsuperscript{2})} & \makecell{지목} & \makecell{이용\newline{}상황} & \makecell{용도지역} & \makecell{도로\newline{}교통} & \makecell{형상\newline{}지세} & \makecell{공시\newline{}기준일} & \makecell{공시지가\newline{}(원/m\textsuperscript{2})} & \makecell{비고} \\
\hline
\multirow{2}{*}{\makecell{①}} & \multirow{2}{*}{\makecell{D동\newline{}121}} & \multirow{2}{*}{\makecell{1,452.0}} & \multirow{2}{*}{\makecell{전}} & \multirow{2}{*}{\makecell{전}} & \multirow{2}{*}{\makecell{자연녹지}} & \multirow{2}{*}{\makecell{맹지}} & \multirow{2}{*}{\makecell{부정형\newline{}완경사}} & \makecell{2018.01.01} & \makecell{156,000} & \multirow{2}{*}{\makecell{사업\newline{}구역 내\newline{}(공원\newline{}100\%)}} \\
\cline{9-10}
 &  &  &  &  &  &  &  & \makecell{2019.01.01} & \makecell{160,000} &  \\
\hline
\multirow{2}{*}{\makecell{②}} & \multirow{2}{*}{\makecell{D동\newline{}214-1}} & \multirow{2}{*}{\makecell{2,564.0}} & \multirow{2}{*}{\makecell{과수원}} & \multirow{2}{*}{\makecell{과수원}} & \multirow{2}{*}{\makecell{자연녹지}} & \multirow{2}{*}{\makecell{세로\newline{}(불)}} & \multirow{2}{*}{\makecell{부정형\newline{}완경사}} & \makecell{2018.01.01} & \makecell{166,000} & \multirow{2}{*}{\makecell{사업\newline{}구역 내\newline{}(공원\newline{}100\%)}} \\
\cline{9-10}
 &  &  &  &  &  &  &  & \makecell{2019.01.01} & \makecell{171,000} &  \\
\hline
\multirow{2}{*}{\makecell{③}} & \multirow{2}{*}{\makecell{D동\newline{}산72-4}} & \multirow{2}{*}{\makecell{4,028.0}} & \multirow{2}{*}{\makecell{임야}} & \multirow{2}{*}{\makecell{자연림}} & \multirow{2}{*}{\makecell{자연녹지}} & \multirow{2}{*}{\makecell{맹지}} & \multirow{2}{*}{\makecell{부정형\newline{}완경사}} & \makecell{2018.01.01} & \makecell{28,000} & \multirow{2}{*}{\makecell{사업\newline{}구역 내\newline{}(공원\newline{}100\%)}} \\
\cline{9-10}
 &  &  &  &  &  &  &  & \makecell{2019.01.01} & \makecell{29,000} &  \\
\hline
\multirow{2}{*}{\makecell{④}} & \multirow{2}{*}{\makecell{D동\newline{}457}} & \multirow{2}{*}{\makecell{1,321.0}} & \multirow{2}{*}{\makecell{잡종지}} & \multirow{2}{*}{\makecell{전기타}} & \multirow{2}{*}{\makecell{자연녹지}} & \multirow{2}{*}{\makecell{세로\newline{}(불)}} & \multirow{2}{*}{\makecell{부정형\newline{}완경사}} & \makecell{2018.01.01} & \makecell{260,000} & \multirow{2}{*}{\makecell{사업\newline{}구역 외}} \\
\cline{9-10}
 &  &  &  &  &  &  &  & \makecell{2019.01.01} & \makecell{290,000} &  \\
\hline
\end{tabular}\par}

※ 본 사업구역 내에 소재하는 공시지가 표준지는 모두 3개로, 표준지 기호 ①과 ③은 1단계 사업지 내에 그리고 기호 ②는 2단계 사업지 내에 각각 소재함 \par
※ 2018. 01. 01∼2019. 01. 01 기간 중 B도 C시의 표준지 공시지가 평균변동률은 7.216 \%임 \par
2. 매매사례 \par
{\small\centering\footnotesize\setlength{\tabcolsep}{3pt}
\begin{tabular}{|c|c|c|c|c|c|c|c|c|}
\hline
\makecell{기호} & \makecell{소재지} & \makecell{거래일자} & \makecell{지목} & \makecell{용도지역} & \makecell{면적\newline{}(m\textsuperscript{2})} & \makecell{이용\newline{}상황} & \makecell{거래가액\newline{}(원/m\textsuperscript{2})} & \makecell{비고} \\
\hline
\makecell{가} & \makecell{D동 137} & \makecell{2018. 09. 01} & \makecell{전} & \makecell{자연녹지} & \makecell{1,208.0} & \makecell{전} & \makecell{280,000} & \makecell{*} \\
\hline
\makecell{나} & \makecell{D동 648} & \makecell{2018. 01. 06} & \makecell{전} & \makecell{자연녹지} & \makecell{1,532.0} & \makecell{전기타} & \makecell{360,000} & \makecell{**} \\
\hline
\end{tabular}\par}

* 기호 가: 사업구역 내(공원 100 \%) 토지로, 본건 토지보다 제반 개별요인 우세함 \par
** 기호 나: 사업구역 외 토지로, 인근의 매매가격 수준 및 평가사례 등에 비추어 정상적인 매매로서 당해 사업으로 인한 영향을 받지 아니한 것으로 판단됨 \par
3. 평가사례 \par
{\small\centering\footnotesize\setlength{\tabcolsep}{3pt}
\begin{tabular}{|c|c|c|c|c|c|c|c|c|}
\hline
\makecell{기호} & \makecell{소재지} & \makecell{기준시점} & \makecell{평가목적\newline{}(사업명칭)} & \makecell{지목} & \makecell{면적\newline{}(m\textsuperscript{2})} & \makecell{용도지역} & \makecell{평가액\newline{}(원/m\textsuperscript{2})} & \makecell{비고} \\
\hline
\makecell{ㄱ} & \makecell{A동\newline{}1207} & \makecell{2019. 02. 01} & \makecell{보상\newline{}(△△테마공원\newline{}주차장조성사업)} & \makecell{전} & \makecell{2,004.0} & \makecell{자연녹지} & \makecell{320,000} & \makecell{*} \\
\hline
\makecell{ㄴ} & \makecell{E동\newline{}36} & \makecell{2018. 04. 08} & \makecell{보상\newline{}(중로 3-XX호\newline{}개설공사)} & \makecell{전} & \makecell{1,082.0} & \makecell{자연녹지} & \makecell{380,000} & \makecell{**} \\
\hline
\end{tabular}\par}

* 기호 ㄱ: 전체 65필지 중 협의체결률은 약 45 \%로서, 가격조사일 현재 나머지 필지는 수용재결 절차에 있는 것으로 조사됨 \par
** 기호 ㄴ: 본 사업구역이 소재하는 D동과 인근의 E동을 남북으로 연결하는 도로 사업임 \par
< 자료 4 > 지가변동률 등 \par
{\small\centering
\begin{tabular}{|c|c|c|c|}
\hline
\multirow{2}{*}{\makecell{기간}} & \multicolumn{2}{|c|}{\makecell{변동률(\%)}} & \multirow{2}{*}{\makecell{비고}} \\
\cline{2-3}
 & \makecell{평균} & \makecell{녹지} &  \\
\hline
\makecell{2018. 01. 01∼2019. 06. 29} & \makecell{4.108} & \makecell{4.202} & \makecell{C시} \\
\hline
\makecell{2019. 01. 01∼2019. 06. 29} & \makecell{1.403} & \makecell{1.470} & \makecell{C시} \\
\hline
\makecell{2018. 10. 01∼2019. 06. 29} & \makecell{2.567} & \makecell{2.718} & \makecell{C시} \\
\hline
\makecell{2018. 10. 01∼2019. 06. 29} & \makecell{2.479} & \makecell{2.692} & \makecell{B도} \\
\hline
\makecell{2018. 01. 06∼2019. 06. 29} & \makecell{3.549} & \makecell{3.892} & \makecell{C시} \\
\hline
\makecell{2018. 04. 08∼2019. 06. 29} & \makecell{3.002} & \makecell{3.112} & \makecell{C시} \\
\hline
\makecell{2018. 09. 01∼2019. 06. 29} & \makecell{2.651} & \makecell{2.847} & \makecell{C시} \\
\hline
\makecell{2019. 02. 01∼2019. 06. 29} & \makecell{1.082} & \makecell{1.236} & \makecell{C시} \\
\hline
\end{tabular}\par}

※ 생산자물가상승률은 인근지역의 적정한 지가변동을 반영하고 있지 아니하다고 판단하여 검토 생략 \par
< 자료 5 > 요인비교 자료 \par
1. 지역요인: 본건 및 공시지가 표준지, 매매사례 및 평가사례는 모두 인근지역에 소재하여 지역요인 대등함 \par
2. 개별요인 \par
(공시지가 표준지: 1.00) \par
{\small\centering
\begin{tabular}{|c|c|c|c|c|c|}
\hline
\multirow{2}{*}{\makecell{공시지가\newline{}표준지}} & \multirow{2}{*}{\makecell{본건\newline{}(연번 1)}} & \multicolumn{2}{|c|}{\makecell{매매사례}} & \multicolumn{2}{|c|}{\makecell{평가사례}} \\
\cline{3-6}
 &  & \makecell{가} & \makecell{나} & \makecell{ㄱ} & \makecell{ㄴ} \\
\hline
\makecell{①} & \makecell{1.00} & \makecell{1.04} & \makecell{1.08} & \makecell{1.00} & \makecell{1.12} \\
\hline
\makecell{②} & \makecell{0.95} & \makecell{0.98} & \makecell{0.95} & \makecell{0.95} & \makecell{0.95} \\
\hline
\makecell{③} & \makecell{4.00} & \makecell{4.10} & \makecell{4.05} & \makecell{4.00} & \makecell{4.00} \\
\hline
\makecell{④} & \makecell{0.90} & \makecell{0.92} & \makecell{0.97} & \makecell{0.90} & \makecell{1.00} \\
\hline
\end{tabular}\par}

※ 상기 개별요인 비교 자료는 도시계획시설 공원 저촉에 따른 제한을 반영하지 않은 수치이며 인근지역에 대한 매매사례 기타 평가사례 등에 대한 조사 결과, 도시계획시설 공원에 저촉된 것에 비해 저촉되지 아니한 상태로의 가치 상승률은 20 \%(임야) ∼ 80 \%(대) 수준을 나타내고 있는바, 필요 시 공원 저촉 여부에 따른 추가 요인비교를 행함 \par
\end{longquestion}
